%----------------------------------------------------------------------------------------
% NOTATIE EN FIGUREN VAN DE CURSUS FINANCIËLE ALGEBRA
% Staat in de preambule en niet in een hoofdstuk: ook de lijsten met oplossingen gebruiken dit.
%----------------------------------------------------------------------------------------
% Bewaarde grafiek, geplaatst door maak-pdf.py: bestand en pagina in figuren-cache/
\newcommand{\grafiekbestand}[2]{\includegraphics[page=#2,trim=3bp 3bp 3bp 3bp]{figuren-cache/#1}}

\providecommand{\N}{\mathbb{N}}
\providecommand{\R}{\mathbb{R}}
\newcommand{\dx}{\,\mathrm{d}x}
\newcommand{\dt}{\,\mathrm{d}t}
\newcommand{\e}{\mathrm{e}}
\newcommand{\ig}{\displaystyle\int}
\newcommand{\grens}[3]{\Bigl[#1\Bigr]_{#2}^{#3}}
% Rechter- en linkerlimiet, met de ongelijkheid onder de pijl
\newcommand{\limr}[2]{\lim\limits_{\substack{#1 \to #2 \\ >}}}
\newcommand{\liml}[2]{\lim\limits_{\substack{#1 \to #2 \\ <}}}

% Zoals stappen, maar met een bredere kolom voor de berekening: voor lange formules
\newenvironment{stappenbreed}{%
  \par\renewcommand{\arraystretch}{1.25}%
  \noindent\begin{tabular}{@{}>{$\displaystyle}p{9.8cm}<{$}@{\hspace{4mm}}>{\RaggedRight\small\itshape\arraybackslash}p{4.9cm}@{}}%
}{\end{tabular}\par}

% Zoals stappen, met eigen kolombreedtes: \begin{stappenvrij}{6.2cm}{3.8cm}. Voor naast een figuur.
\newenvironment{stappenvrij}[2]{%
  \par\renewcommand{\arraystretch}{1.25}%
  \noindent\begin{tabular}{@{}>{$\displaystyle}p{#1}<{$}@{\hspace{3mm}}>{\RaggedRight\small\itshape\arraybackslash}p{#2}@{}}%
}{\end{tabular}\par}

% Tekst of berekening links, figuur rechts: \naast[breedte rechts]{links}{rechts}
% De breedte rechts komt uit de figuur zelf; het optionele argument wordt niet gebruikt.
\newsavebox{\naastbox}
\newcommand{\naast}[3][]{%
  \par\noindent\sbox{\naastbox}{#3}%
  \begin{minipage}[t]{\dimexpr\linewidth-\wd\naastbox-3mm\relax}\vspace{0pt}#2\end{minipage}\hfill
  \begin{minipage}[t]{\wd\naastbox}\vspace{0pt}\usebox{\naastbox}\end{minipage}\par\addvspace{1mm}}
% Een rij figuren, gecentreerd
\newcommand{\figuren}[1]{\par\addvspace{1mm}{\centering #1\par}\addvspace{1.5mm}}
% Opgave met een grafiek rechts: \begin{metgrafiek}[breedte]{tikz} opgave en deelvragen \end{metgrafiek}
\newsavebox{\grafiekbox}
\newenvironment{metgrafiek}[2][]{%
  \sbox{\grafiekbox}{#2}%
  \begin{minipage}[t]{\dimexpr\linewidth-\wd\grafiekbox-3mm\relax}%
}{%
  \end{minipage}\hfill
  \begin{minipage}[t]{\wd\grafiekbox}\vspace{-2mm}\usebox{\grafiekbox}\end{minipage}\par}

% Tabel met een raster, alles gecentreerd: \begin{rastertabel}{ccc} ... \end{rastertabel}
\newenvironment{rastertabel}[1]{\par\addvspace{1mm}\begingroup\centering\renewcommand{\arraystretch}{1.3}%
  \setlength{\arrayrulewidth}{.4pt}\arrayrulecolor{lijngrijs}\begin{tabular}{#1}}%
  {\end{tabular}\par\endgroup\addvspace{1.5mm}}

% Rooster met assen: \rooster{xmin}{xstap}{xmax}{ymin}{ystap}{ymax}, met getallen op de assen (gehele waarden).
% \roosterleeg doet hetzelfde zonder getallen, voor eigen labels.
\newcommand{\roosterleeg}[6]{%
  \draw[black!14,very thin,xstep=#2,ystep=#5] (#1,#4) grid (#3,#6);
  \draw[->,grijs] (#1,0) -- ({#3+.4*#2},0) node[right,text=black] {$x$};
  \draw[->,grijs] (0,#4) -- (0,{#6+.4*#5}) node[above,text=black] {$y$};}
\newcommand{\rooster}[6]{%
  \roosterleeg{#1}{#2}{#3}{#4}{#5}{#6}
  \pgfmathtruncatemacro{\rtweex}{#1+#2}\pgfmathtruncatemacro{\rtweey}{#4+#5}
  \foreach \t in {#1,\rtweex,...,#3}{\ifnum\t=0\else\node[below,inner sep=1.5pt,text=black] at (\t,0) {$\t$};\fi}
  \foreach \t in {#4,\rtweey,...,#6}{\ifnum\t=0\else\node[left,inner sep=1.5pt,text=black] at (0,\t) {$\t$};\fi}}
% Assen zonder rooster
\newcommand{\assen}[4]{%
  \draw[->,grijs] (#1,0) -- (#2,0) node[right,text=black] {$x$};
  \draw[->,grijs] (0,#3) -- (0,#4) node[above,text=black] {$y$};}
% In de functievoorschriften is \x de veranderlijke; schrijf (\x)^2 en sin(\x r).
\newcommand{\kromme}[4][]{\draw[accent,thick,#1] plot[domain=#2:#3,samples=90,smooth] (\x,{#4});}
% \gebied: tussen de grafiek en de horizontale as
\newcommand{\gebied}[4][accent!30]{\fill[#1] plot[domain=#2:#3,samples=70] (\x,{#4}) -- (#3,0) -- (#2,0) -- cycle;}
% \riemann{aantal}{a}{b}{voorschrift}{plaats}: rechthoeken, met de hoogte op plaats 0 (links), 0.5 (midden) of 1 (rechts)

%----------------------------------------------------------------------------------------
% Eigen voor Financiële algebra
%----------------------------------------------------------------------------------------
% Rentevoet per maand, trimester, semester: i_{12}, i_4, i_2; rentefactor u = 1 + i
\newcommand{\im}{i_{12}}
\newcommand{\um}{u_{12}}
% Een bedrag in euro in een formule: \bedrag{1\,234{,}56}
\newcommand{\bedrag}[1]{#1\ \text{euro}}

% Tijdlijn: \tijdas{lengte}{lijst positie/label}: een pijl met streepjes en labels eronder (in een tikzpicture)
\newcommand{\tijdas}[2]{%
  \draw[->,grijs] (-.25,0) -- ({#1+.45},0);
  \foreach \t/\l in {#2}{\draw[grijs] (\t,.07) -- (\t,-.07) node[below,text=black,inner sep=2pt] {\l};}}
% Een storting of betaling op de tijdlijn: \betaling{positie}{label boven}
\newcommand{\betaling}[2]{\fill[accent] (#1,0) circle (1.7pt) node[above=1pt,text=black] {#2};}

% Rekenblad: \begin{rekenblad}{aantal kolommen}{kolomletters A & B & ...} rijen met \rij{nummer} & ... \\ \end{rekenblad}
\newenvironment{rekenblad}[2]{\par\addvspace{1mm}\begingroup\centering\small\renewcommand{\arraystretch}{1.25}%
  \setlength{\arrayrulewidth}{.4pt}\arrayrulecolor{lijngrijs}\setlength{\tabcolsep}{4pt}%
  \begin{tabular}{|>{\columncolor{black!8}}c|*{#1}{r|}}\hline
  \rowcolor{black!8} & #2 \\ \hline}%
  {\end{tabular}\par\endgroup\addvspace{1.5mm}}
